% ============================================================
%  文档类选项（在 main.tex 的 \documentclass[...]{tongjithesis} 中设置）
%  field=science    理工科（工科类、理科类，默认）— 章节编号：1 / 1.1 / 1.1.1
%  field=humanities 文科  （人文、法学、外语、艺术）        — 章节编号：一、/（一）/ 1.
%  经管类属社科：编号要求依理工，撰写依理工模版；26 届为过渡期，可选 science 或
%  humanities；自 27 届起统一使用 science。
% ============================================================

% ============================================================
%  封面信息
% ============================================================
\school{计算机科学与技术学院}
\major{计算机科学与技术}
\student{2251651}{郭冰放}
\thesistitle{基于RISC-V的UnixV6++操作系统核心框架与进程管理移植}{}
\thesistitleeng{Porting the Core Framework and Process Management of the UnixV6++ Operating System Based on RISC-V}{}
\thesisadvisor{方钰\quad 教授}
\thesisdate{\the\year}{\the\month}{\the\day} % 使用当前日期; 也可以自定义日期

% ============================================================
%  信息说明页
% ============================================================

% 成果类型（三选一）：
%   thesis      — 毕业论文，摘要
%   design      — 毕业设计（软件/创意/建筑等非工程设计类），摘要
%   engineering — 毕业设计（工程设计类），设计总说明
\infotype{thesis}

% 内容简述（请用300字以内简要概述）
\infoabstract{%
本文研究 Unix V6++ 教学操作系统面向 RISC-V 平台的移植与 Rust 重构。重构前系统以 IA-32/C++/PC 设备模型为基础，当前系统运行于 QEMU virt 与 OpenSBI 环境，内核主体采用 Rust 实现，用户态 C 程序交叉编译为 ELF64/RISC-V 可执行文件。论文压缩为五章，重点围绕 RISC-V 启动与 Trap、系统调用、进程管理、ELF 加载、信号、重构差异、资源生命周期、TTY/Shell 交互和 QEMU 测试展开，文件系统和内存管理主要作为运行支撑条件处理。系统已覆盖构建、镜像生成和交互验证流程，并保留 Unix V6++ 的主要教学语义。
}

% 毕业设计：图纸数量与正文字数（成果类型为 design 或 engineering 时填写，两者须同时填写）
% \infodrawings{5}
% \infowordcount{15000}

% 毕业论文：正文总字数（成果类型为 thesis 时填写）
\infothesiswords{17800}

% 随附资料（每条用 \item 开头；不填则显示空白行）
\infomaterials{
  \item 程序源代码；
  \item 构建脚本与测试用用户程序；
  \item 论文插图源文件。
}
